\chapter*{\begin{center}
\vspace{-2.5em}	РЕФЕРАТ \vspace{-1.5em}
\end{center}
}
\addcontentsline{toc}{chapter}{РЕФЕРАТ}

Расчётно-пояснительная записка 33 страницы, 15 источников, 1 приложение.

МЕТОДЫ СИНХРОНИЗАЦИИ, СЕТЕВОЕ ВЗАИМОДЕЙСТВИЕ, ОБРАБОТКА ИГРОВЫХ СОБЫТИЙ.

Объектом исследования стали методы синхронизации состояния в многопользовательских игровых системах.

Цель работы -- сравнительный анализ методов предотвращения рассинхронизации в многопользовательских
игровых приложениях.

В работе выполнен анализ предметной области синхронизации игровых состояний, исследованы существующие 
методы синхронизации данных, сформулированы критерии их оценки и проведён сравнительный анализ на основе 
выделенных параметров.